\chapter{Wstęp}
\section{Kontekst i tło problemu}
Współczesne systemy informatyczne opierają się w przeważającej mierze na architekturze rozproszonych mikro-usług, których orkiestracja odbywa się przy pomocy platform takich jak Kubernetes lub innych wewnętrznych rozwiązań. Skala tych środowisk nieustannie rośnie – dobitnym przykładem są globalne przedsiębiorstwa technologiczne, takie jak np. Uber, który operuje na zawiłej sieci ok. 2200 krytycznych mikro-usług\cite{uber_microservices}, podobnie jak Spotify \cite{spotify_microservices}. Choć podejście oparte na konteneryzacji przyniosło ogromne korzyści w zakresie skalowalności, elastyczności i niezależności wdrożeń, drastycznie zwiększyło ono również złożoność operacyjną środowisk IT. Zapewnienie ciągłej operacyjności (High Availability) tak skomplikowanych systemów rozproszonych pochłania gigantyczne nakłady finansowe. Koszty te wynikają w głównej mierze z konieczności opłacania i utrzymywania wysoce wykwalifikowanych zespołów inżynierskich (Site Reliability Engineering / DevOps), które zmuszone są do pracy w systemie całodobowych dyżurów (tzw. "on-call"). Kiedy dochodzi do awarii, presja czasu jest ogromna. Diagnoza i przywrócenie systemu do stanu używalności (remediacja) wymagają od dyżurujących inżynierów analizy setek powiązań sieciowych, rozproszonych logów, metryk i konfiguracji klastra. Stanowi to nie tylko rosnące obciążenie dla zespołów IT, ale również realne straty biznesowe z powodu przedłużających się przestojów. Dla przykładu platforma Amazon Prime raportuje, że każda godzina przestoju kosztuje firmę 100 milionów dolarów \cite{amazon_prime_down_cost}. Kluczowym rozwiązaniem tego narastającego problemu staje się zaawansowana automatyzacja procesów naprawczych. Jak zauważają autorzy pracy "Leveraging Large Language Models for the Auto-remediation of Microservice Applications: An Experimental Study" \cite{sarda2024leveraging}, Duże Modele Językowe (LLM) wykazują ogromny potencjał w automatyzacji zadań związanych z diagnozowaniem i rozwiązywaniem incydentów w mikro-usługach. Jednakże standardowe, jednoagentowe podejście nie jest w stanie efektywnie wspierać remediacji w środowisku o wysokiej złożoności, w którym niezbędna jest rozległa sieć wyspecjalizowanych narzędzi i wiedzy domenowej.
\section{Cel badawczy i wykorzystanie architektury wieloagentowej}
Aby przezwyciężyć ograniczenia pojedynczych modeli, coraz częściej sięga się po systemy wieloagentowe (Multi-Agent Systems – MAS). 
Wdrożenie takich systemów wiąże się jednak z licznymi wyzwaniami dotyczącymi koordynacji, komunikacji i stabilności działania \cite{han2024llm_multiagent_challenges}. 
W krytycznych zastosowaniach, takich jak zarządzanie infrastrukturą, kluczowe jest zapewnienie rzetelności podejmowanych działań, co podkreślono w badaniach w Industrial A.I. Laboratory Hitachi America, Ltd. \cite{lee2025reliable_multiagent_llm}. Głównym celem niniejszej pracy jest eksperymentalne zbadanie, w jaki sposób system wieloagentowy może adaptować się do zmieniających się warunków oraz zwiększać skuteczność i niezawodność autoremediacji aplikacji mikrousługowej. Aby zapewnić ciągłe uczenie się systemu i poprawę jakości wdrożeń, w pracy zaimplementowano potok inspirowany publikacją \textit{Agentic Context Engineering: Evolving Contexts for Self-Improving Language Models} \cite{zhang2025agentic_context_engineering} zaadaptowany do bardziej złożonej architektury systemu agentowego. Metoda ta pozwala agentom na dynamiczne tworzenie i modyfikowanie własnej bazy wiedzy (tzw. Playbook lub SOP) w czasie rzeczywistym.

\section{Teza i pytania badawcze}
Główna teza pracy zakłada, że zastosowanie wieloagentowego systemu wyposażonego w adaptacyjny potok ewolucji kontekstu (\textit{Agentic Context Engineering}) istotnie podnosi skuteczność autoremediacji aplikacji mikrousługowej --- mierzoną wskaźnikami jakościowymi (poprawność diagnozy, skuteczność naprawy, widoczna regeneracja systemu w metrykach) --- przy akceptowalnym koszcie operacyjnym wyrażonym w czasie do rozwiązania incydentu (TTR, \textit{Time To Resolution}) oraz liczbie iteracji agentów.

W celu weryfikacji postawionej tezy sformułowano następujące pytania badawcze:

1. W jaki sposób dynamiczna ewolucja instrukcji operacyjnych (Playbooka) wpływa na skuteczność agentów w obliczu różnorodnych klas incydentów --- od problemów infrastrukturalnych, przez błędy konfiguracyjne, po regresje w kodzie aplikacji?

2. Jakim kosztem operacyjnym --- mierzonym liczbą iteracji, czasem do rozwiązania incydentu (TTR), intensywnością wywołań narzędzi oraz odsetkiem błędnych wywołań --- okupiony jest wzrost skuteczności remediacji wynikający z ewolucji Playbooka?

3. Jakie emergentne wzorce organizacji pracy systemu wieloagentowego (w szczególności w zakresie głębokości drzewa zadań i hierarchicznej delegacji pomiędzy wyspecjalizowanymi agentami) pojawiają się w wyniku adaptacji kontekstu w pętli ACE?

4. Czy hierarchiczna propagacja refleksji oraz mechanizm \textit{Reflektora} i \textit{Kuratora} w pętli ACE skutecznie ograniczają regresje wprowadzane przez nowo wygenerowane, błędne lub nadmiernie ogólne reguły naprawcze?

\section{Zakres i struktura pracy}

Zakres pracy obejmuje: (1)~opracowanie hierarchicznej architektury wieloagentowej w paradygmacie \textit{agent-as-a-tool}, (2)~implementację izolowanego, deterministycznego środowiska badawczego opartego na klastrze Kubernetes, aplikacji mikrousługowej \textit{Robot Shop} oraz autorskim komponencie wstrzykującym usterki (\textit{Saboteur Agent}), (3)~adaptację algorytmu \textit{Agentic Context Engineering} do kontekstu wieloagentowego, w szczególności poprzez hierarchiczną propagację refleksji wzdłuż drzewa zadań oraz mechanizmy ograniczania stronniczości reguł Playbooka. 

Ponadto zakres obejmuje (4)~przeprowadzenie i ilościową analizę eksperymentu porównawczego pomiędzy wariantem referencyjnym (bez uczenia) a wariantem z aktywnym potokiem ACE, na deterministycznej sekwencji 45 incydentów reprezentujących cztery klasy usterek. Poniżej zestawiono treść kolejnych rozdziałów.

\textit{Rozdział II} opisuje autorskie środowisko testowe oparte na aplikacji \textit{Robot Shop} \cite{instana_robot_shop}, w literaturze uznawanej za sensowny punkt odniesienia do ewaluacji systemów mikrousługowych \cite{fischer2024microservice_systems_mapping}, definiuje cztery klasy generowanych incydentów (od hybrydowych usterek \textit{Chaos Mesh}~+~sabotaż kodu po wycieki zasobów) oraz wprowadza zestaw binarnych wskaźników jakościowych (\texttt{root\_cause\_analysis}, \texttt{successful\_fix}, \texttt{system\_recovery\_visible}) wykorzystywanych do automatycznej oceny remediacji. Wskazano także istotne ograniczenia i podatności przyjętego układu.

\textit{Rozdział III} zawiera przegląd literatury dotyczącej agentowych systemów opartych na dużych modelach językowych (paradygmat Re-Act, problem \textit{Lost in the Middle}, kompresja kontekstu) oraz przedstawia projektowanie autorskiego, czteroskładnikowego agenta SRE (\textit{Incident Commander}, \textit{Monitoring Agent}, \textit{DevOps Agent}, \textit{Coder Agent}) wraz z hierarchicznym modelem koordynacji opartym na drzewie zadań (\textit{tasks tree}) i mechanizmach delegacji.

\textit{Rozdział IV} prezentuje szczegółowy opis metody uczenia oraz autorską implementację potoku \textit{Agentic Context Engineering} (ACE) w architekturze wieloagentowej, ze szczególnym uwzględnieniem ról agenta \textit{Reflektora} i \textit{Kuratora}, struktury Playbooka (reguły IF/THEN, parametry \texttt{helpful}/\texttt{harmful}), mechanizmu hierarchicznej propagacji refleksji wzdłuż drzewa zadań oraz technik ograniczania stronniczości (ukrywania wag przed agentem wykonawczym i Reflektorem).

\textit{Rozdział V} przedstawia projekt eksperymentu referencyjnego (\textit{Baseline}, bez aktywnego potoku ACE) oraz konfigurację eksperymentu ewolucyjnego (z aktywnym potokiem, w trybie \textit{batch learning} co 5 epizodów), opisuje rolę agenta sędziego (\textit{Judge Agent}) oraz formalizuje metryki ewaluacji wykorzystywane w analizie wyników.

\textit{Rozdział VI} prezentuje wyniki eksperymentu na deterministycznej sekwencji 45 incydentów: globalne wskaźniki sukcesu, dezagregację wyników w podziale na cztery klasy usterek, chronologiczną dynamikę procesu uczenia, metryki operacyjne (TTR, liczba iteracji, błędy wywołań narzędzi) oraz zjawiska emergentne w organizacji pracy systemu wieloagentowego --- w szczególności pogłębienie struktury drzewa zadań i samoorganizację \textit{Monitoring Agenta} jako huba orkiestrującego wieloetapową analizę obserwowalności.

\textit{Rozdział VII} podsumowuje ustalenia, weryfikuje postawioną tezę oraz wskazuje możliwe kierunki dalszych badań, w tym warunkowanie reguł Playbooka klasą wykrywanego incydentu jako odpowiedź na zaobserwowaną w Klasie~1 regresję.
